% ============================================================================
%  SECCIÓN 1.10: TIPO Y ESTUDIO DE LA INVESTIGACIÓN
% ============================================================================

\section{Tipo y estudio de la investigaci\'on}

En funci\'on al nivel de profundidad que se prev\'e alcanzar, la investigaci\'on
es de tipo exploratoria y descriptiva, seg\'un la
clasificaci\'on propuesta por Dankhe (1986), citada por \citet{hernandez2014}:

\begin{itemize}
  \item Exploratoria: porque la aplicaci\'on de sistemas de
    visualizaci\'on interactiva al sector deportivo nacional es un tema que no
    ha sido estudiado ampliamente en el medio, por lo que la investigaci\'on
    permite familiarizarse con esta realidad poco conocida.

  \item Descriptiva: porque el objetivo consiste en especificar las
    caracter\'isticas y rasgos importantes de los escenarios deportivos a nivel
    nacional, as\'i como las propiedades de la difusi\'on de la oferta
    deportiva, describiendo la situaci\'on actual del objeto de estudio.
\end{itemize}

\begin{table}[H]
  \centering
  \caption{Alcances de la investigaci\'on}
  \label{tab:alcances}
  \begin{tablaAPA}
    \begin{tabular}{@{}p{4.6cm}p{11.2cm}@{}}
      \toprule
      \textbf{Tipo de investigaci\'on} & \textbf{Alcance en el proyecto} \\
      \midrule
      Exploratoria &
      Reconocimiento inicial del sector deportivo y de las soluciones
      tecnol\'ogicas aplicables a la visualizaci\'on interactiva \\
      \addlinespace
      Descriptiva &
      Caracterizaci\'on de los escenarios deportivos y del proceso de difusi\'on
      de la oferta deportiva \\
      \bottomrule
    \end{tabular}
  \end{tablaAPA}
\end{table}

\textbf{FUENTE:} Elaboraci\'on propia en base a \citet{hernandez2014}.
